\documentclass{article}
\usepackage[UTF8]{ctex}
\usepackage{amsmath,amssymb,amsfonts}
\usepackage{enumitem}
\usepackage{geometry}
\geometry{a4paper, margin=1in}
% === 错题本排版优化（不参与题目标识） ===
\setlength{\parindent}{0pt}
\setlength{\parskip}{0.35em plus 0.10em minus 0.08em}
\setlength{\emergencystretch}{3em}
\setlist{itemsep=0.25em, topsep=0.40em, parsep=0pt}
\allowdisplaybreaks[2]
% === 错题本排版优化结束 ===
\begin{document}
\section*{哈尔滨工业大学2025年高等代数真题}
\section*{75．哈尔滨工业大学2025年高等代数试卷}

\subsection{题目1}
1．设 \(A\) 是 \(m \times n\) 矩阵，证明：秩 \((A)=r\) 的充分必要条件是存在 \(m \times r\) 矩阵 \(B\) 和 \(r \times n\) 矩阵 \(C\) ，满足秩 \((B)=\) 秩 \((C)=r\) ，且 \(A=B C\) 。

\subsection{题目2}
2．计算行列式

\[
\left|\begin{array}{cccc}
a_{1}^{2} & a_{1} a_{2}+1 & \cdots & a_{1} a_{n}+1 \\
a_{2} a_{1}+1 & a_{2}^{2} & \cdots & a_{2} a_{n}+1 \\
\vdots & \vdots & & \vdots \\
a_{n} a_{1}+1 & a_{n} a_{2}+1 & \cdots & a_{n}^{2}
\end{array}\right|
\]

\subsection{题目3}
3．设 \(p(x)\) 是数域 \(F\) 上次数大于等于 1 的多项式。证明下列说法等价：

\begin{itemize}
  \item[（1）] \(p(x)\) 是不可约多项式．
  \item[（2）] 对任意的 \(f(x), g(x) \in F[x]\) ，如果 \(p(x) \mid f(x) g(x)\) ，则有 \(p(x) \mid f(x)\) 或 \(p(x) \mid g(x)\).
  \item[（3）] 对任意的 \(f(x) \in F[x]\) ，都有 \(p(x) \mid f(x)\) 或 \((p(x), f(x))=1\) ．
\end{itemize}

\subsection{题目4}
4．讨论当 \(\lambda\) 取何值时，如下方程有解，并求解．

\[
\left\{\begin{array}{l}
x_{1}+\lambda x_{2}+\lambda x_{3}+x_{4}=0 \\
2 x_{1}+x_{2}+x_{3}+2 x_{4}=0 \\
3 x_{1}+(\lambda+2) x_{2}+(\lambda+4) x_{3}+4 x_{4}=1
\end{array}\right.
\]

\subsection{题目5}
5．设 \(A\) 是 \(m \times n\) 矩阵，证明：非齐次线性方程组 \(A X=\beta\) 有解的充分必要条件是齐次线性方程组\(A^{\mathrm{T}} Y=0\) 的任意解 \(\alpha\) 满足 \(\alpha^{\mathrm{T}} \beta=0\) ．

\subsection{题目6}
6．设实数 \(x_{1}, x_{2}, x_{3}, x_{4}\) 满足 \(x_{1}^{2}+x_{2}^{2}+x_{3}^{2}+x_{4}^{2}=1\) ，求

\[
y=x_{1} x_{2}+x_{2} x_{3}+x_{3} x_{4}+x_{4} x_{1}-2\left(x_{1} x_{3}+x_{2} x_{4}\right)
\]

的取值范围．

\begin{itemize}

\subsection{题目7}
\end{itemize}

\subsection{题目7}
原题边界“7．”：设实数序列 \(x_{0}=1, x_{1}=12, x_{2}=-10, x_{3}, x_{4}, \cdots, x_{n}, \cdots\) 满足
\[
x_{n+3}+30 x_{n}=31 x_{n+1}, n \geq 0 .
\]
求 \(\left\{x_{n}\right\}\) 的通项公式。

\subsection{题目8}

\subsection{题目8}
原题边界“8．”：设矩阵 \(A=\left(\begin{array}{ll}-5 & 5 \\ -6 & 6\end{array}\right)\) ，求满足 \(A B+B A=A B A\) 的所有复矩阵 \(B\) 。

\subsection{题目9}
9．设矩阵 \(A\) 的特征多项式为 \(\lambda^{3}-6 \lambda^{2}+32\) ，求 \(A\) 的若尔当标准形．

\subsection{题目10}
10．设 \(a\) 为实数，矩阵

\[
A=\left(\begin{array}{ll}
0 & a \\
1 & 1
\end{array}\right)
\]

若存在正交矩阵 \(P\) ，使得 \(P^{-1} A P\) 为上三角的实矩阵，求 \(a\) 的取值范围．
\end{document}
